% Generated draft: requires mathematical and source review.
\documentclass[11pt]{article}
\usepackage[margin=25mm]{geometry}
\usepackage{fontspec}
\setmainfont{Latin Modern Roman}
\usepackage{amsmath,amssymb,mathtools,unicode-math}
\setmathfont{Latin Modern Math}
\usepackage{enumitem,needspace,xurl,hyperref}
\setlength{\parindent}{0pt}
\setlength{\parskip}{5pt}
\setlength{\emergencystretch}{3em}
\newcommand{\N}{\mathbb{N}}
\newcommand{\Z}{\mathbb{Z}}
\newcommand{\Q}{\mathbb{Q}}
\newcommand{\R}{\mathbb{R}}
\newcommand{\C}{\mathbb{C}}
\newcommand{\F}{\mathbb{F}}
\newcommand{\A}{\mathbb{A}}
\newcommand{\PP}{\mathbb{P}}
\newcommand{\E}{\mathbb{E}}
\newcommand{\eps}{\varepsilon}
\newcommand{\dd}{\,\mathrm{d}}
\begin{document}
\section*{UnsolvedMath ID: 20000971}
\textbf{Source problem code:} catalogue entry
\subsection{Definitions and mathematical statement}
The source discusses the ``explorers game,'' interpreted by the attached aid as internal diffusion limited aggregation (IDLA): particles start at a fixed origin one at a time, follow independent random walks, and each joins the aggregate at its first visit to an unoccupied vertex. Write $A_n$ for the occupied set after $n$ particles. The precise walk and boundary rule on a restricted lattice are essential parts of any proposed model and are not specified by the workshop question.

Persi Diaconis asks for growth models with fluctuations at scales intermediate between the one-dimensional $\sqrt n$ and two-dimensional $\log n$ scales quoted in the discussion; it also quotes $\sqrt{\log n}$ in dimensions at least three and a ball radius of order $n^{1/d}$. These are historical source descriptions, with no specified probability mode, fluctuation statistic, or centering convention. Andrea Sportiello asks whether a lattice wedge has a circular-sector limit shape and proposes domains $|y|\leq x^{1/\gamma}$ to interpolate dimensions. Charles Smart comments that behavior away from the border should be unaffected. Igor Pak proposes more general domains $|y|\leq f(x)$ to tune fluctuation growth. Establishing a limit shape and a defined fluctuation scale for such models, rather than assuming them, is the target.

The aid treats a specific induced graph. For fixed $0<\alpha<1$ let
\[
 V(G_\alpha)=\{(k,j)\in\Z^2:k\geq0,\ |j|\leq\lfloor k^\alpha\rfloor\},
\]
with nearest-neighbor lattice edges, unit conductances, and origin $o=(0,0)$. Here $\alpha=1/\gamma$ corresponds to $\gamma>1$ in Sportiello's proposal. Let $B_\alpha(R)$ be the ball of radius $R$ in graph distance. In the finite part $0\leq k\leq R$, wire together all vertices of the terminal column $L_R=\{(R,j):|j|\leq\lfloor R^\alpha\rfloor\}$. The aid records, for integer $R\to\infty$ with $\alpha$ fixed,
\[
 |B_\alpha(R)|=\frac{2}{1+\alpha}R^{1+\alpha}+O_\alpha(R+R^{2\alpha}),\qquad
 R_{\mathrm{eff}}(o\leftrightarrow L_R)\sim\frac{R^{1-\alpha}}{2(1-\alpha)}.
\]
Its upper resistance construction is reported to have energy within $O_\alpha(\log R)$ of the Nash--Williams cutset lower sum. These are geometric and electrical-network statements, not IDLA fluctuation estimates.

The aid separately describes a continuum Neumann cone benchmark whose occupied region is a circular sector. For opening angle $\theta$ and unit area density, area $n$ corresponds to radius $\sqrt{2n/\theta}$, so the angle affects mass--radius normalization even when the boundary curve is circular. This deterministic benchmark does not prove that lattice IDLA in a wedge has that shape. In particular neither the volume/resistance asymptotics nor this benchmark establishes any intermediate IDLA fluctuation exponent.
\subsection{Short English statement}
Find growth models with fluctuations between logarithmic and square-root size, especially IDLA in narrowing lattice domains. The supplied aid computes volume and resistance for parabolic induced graphs and a deterministic cone benchmark; it does not prove an IDLA shape or fluctuation law.
\subsection{Sources}
\begin{itemize}
\item AIM workshop \emph{Generalizations of chip-firing and the critical group}, July 8--12, 2013, problem 6, attributed discussion by Persi Diaconis, Andrea Sportiello, Charles Smart, and Igor Pak: \url{https://aimath.org/pastworkshops/chipfiringproblems.pdf}. The supplied verification notes that the list's introduction describes comments as recorder's summaries, not necessarily verbatim quotations.
\item The discussion cites its references [22], [4], and [21] for higher-dimensional results; their bibliographic entries are absent from this record and are not reconstructed.
\item Local registry \texttt{data/unsolvedmath/problems.json}, permanent ID 20000971, exact source code \texttt{catalogue entry}; canonical location \texttt{aim-combinatorics-notes.json notes[95]}. Attached attempt 1, partial result, problem status at run partially solved, rigor proved, low novelty confidence; machine-generated, not peer reviewed, requiring independent verification before citation.
\end{itemize}
\end{document}
